% ============================================================
% content/questions.tex —— 2019 全国Ⅰ卷（理综物理部分）· 试题
% 来源：中国教育在线公开原卷 Word 版（含官方答案）
% 题号保留官方编号 14--21、22--25、33--34；第 1--13 题为理综其他科目。
% 图形见 figures/（全部为 TikZ/circuitikz/pgfplots 矢量重绘）
% ============================================================

\begin{gknotice}
  \textbf{卷种说明：}本卷为 2019 年普通高等学校招生全国统一考试理科综合能力测试（全国卷Ⅰ）的物理部分，
  满分 110 分，考试用时 90 分钟。选择题题号从 14 开始，因为第 1--13 题为理综其他科目试题，
  本工程只收录物理部分。第 33、34 题为选考题（选修 3-3、选修 3-4），考生按原卷要求任选一题作答。
\end{gknotice}

\gkpart{二、选择题}

\textbf{（本题共 8 小题，每小题 6 分，共 48 分。在每小题给出的四个选项中，第 14--18 题只有一项符合题目要求，第 19--21 题有多项符合题目要求。全部选对的得 6 分，选对但不全的得 3 分，有选错的得 0 分。）}

\begin{question}
\gkti{14}氢原子能级示意图如图所示。光子能量在 \SI{1.63}{eV}--\SI{3.10}{eV} 的光为可见光。要使处于基态（$n=1$）的氢原子被激发后可辐射出可见光光子，最少应给氢原子提供的能量为\gkblank

\gkfig{figures/q14-levels}

\begin{gkopts}(4)
  \task A．\SI{12.09}{eV}
  \task B．\SI{10.20}{eV}
  \task C．\SI{1.89}{eV}
  \task D．\SI{1.51}{eV}
\end{gkopts}
\end{question}

\begin{question}
\gkti{15}如图，空间存在一方向水平向右的匀强电场，两个带电小球 $P$ 和 $Q$ 用相同的绝缘细绳悬挂在水平天花板下，两细绳都恰好与天花板垂直，则\gkblank

\gkfig{figures/q15-balls}

\begin{gkopts}(2)
  \task A．$P$ 和 $Q$ 都带正电荷
  \task B．$P$ 和 $Q$ 都带负电荷
  \task C．$P$ 带正电荷，$Q$ 带负电荷
  \task D．$P$ 带负电荷，$Q$ 带正电荷
\end{gkopts}
\end{question}

\begin{question}
\gkti{16}最近，我国为“长征九号”研制的大推力新型火箭发动机联试成功，这标志着我国重型运载火箭的研发取得突破性进展。若某次实验中该发动机向后喷射的气体速度约为 \SI{3}{\kilo\metre\per\second}，产生的推力约为 \SI{4.8e6}{\newton}，则它在 \SI{1}{\second} 时间内喷射的气体质量约为\gkblank

\begin{gkopts}(4)
  \task A．\SI{1.6e2}{\kilogram}
  \task B．\SI{1.6e3}{\kilogram}
  \task C．\SI{1.6e5}{\kilogram}
  \task D．\SI{1.6e6}{\kilogram}
\end{gkopts}
\end{question}

\begin{question}
\gkti{17}如图，等边三角形线框 $LMN$ 由三根相同的导体棒连接而成，固定于匀强磁场中，线框平面与磁感应强度方向垂直，线框顶点 $M$、$N$ 与直流电源两端相接。已知导体棒 $MN$ 受到的安培力大小为 $F$，则线框 $LMN$ 受到的安培力的大小为\gkblank

\gkfig{figures/q17-frame}

\begin{gkopts}(4)
  \task A．$2F$
  \task B．$1.5F$
  \task C．$0.5F$
  \task D．$0$
\end{gkopts}
\end{question}

\begin{question}
\gkti{18}如图，篮球架下的运动员原地垂直起跳扣篮，离地后重心上升的最大高度为 $H$。上升第一个 $\dfrac{H}{4}$ 所用的时间为 $t_{1}$，第四个 $\dfrac{H}{4}$ 所用的时间为 $t_{2}$。不计空气阻力，则 $\dfrac{t_{2}}{t_{1}}$ 满足\gkblank

\gkfig{figures/q18-jump}

\begin{gkopts}(4)
  \task A．$1<\dfrac{t_{2}}{t_{1}}<2$
  \task B．$2<\dfrac{t_{2}}{t_{1}}<3$
  \task C．$3<\dfrac{t_{2}}{t_{1}}<4$
  \task D．$4<\dfrac{t_{2}}{t_{1}}<5$
\end{gkopts}
\end{question}

\begin{question}
\gkti{19}如图，一粗糙斜面固定在地面上，斜面顶端装有一光滑定滑轮。一细绳跨过滑轮，其一端悬挂物块 $N$，另一端与斜面上的物块 $M$ 相连，系统处于静止状态。现用水平向左的拉力缓慢拉动 $N$，直至悬挂 $N$ 的细绳与竖直方向成 $45^{\circ}$。已知 $M$ 始终保持静止，则在此过程中\gkblank

\gkfig{figures/q19-incline}

\begin{gkopts}(1)
  \task A．水平拉力的大小可能保持不变
  \task B．$M$ 所受细绳的拉力大小一定一直增加
  \task C．$M$ 所受斜面的摩擦力大小一定一直增加
  \task D．$M$ 所受斜面的摩擦力大小可能先减小后增加
\end{gkopts}
\end{question}

\begin{question}
\gkti{20}空间存在一方向与纸面垂直、大小随时间变化的匀强磁场，其边界如图（a）中虚线 $MN$ 所示。一硬质细导线的电阻率为 $\rho$、横截面积为 $S$，将该导线做成半径为 $r$ 的圆环固定在纸面内，圆心 $O$ 在 $MN$ 上。$t=0$ 时磁感应强度的方向如图（a）所示，磁感应强度 $B$ 随时间 $t$ 的变化关系如图（b）所示，则在 $t=0$ 到 $t=t_{1}$ 的时间间隔内\gkblank

\gkfig{figures/q20-ab}

\begin{gkopts}(1)
  \task A．圆环所受安培力的方向始终不变
  \task B．圆环中的感应电流始终沿顺时针方向
  \task C．圆环中的感应电流大小为 $\dfrac{BSr}{4\rho}$
  \task D．圆环中的感应电动势大小为 $\dfrac{\pi r^{2}B}{2t_{1}}$
\end{gkopts}
\end{question}

\begin{question}
\gkti{21}在星球 $M$ 上将一轻弹簧竖直固定在水平桌面上，把物体 $P$ 轻放在弹簧上端，$P$ 由静止向下运动，物体的加速度 $a$ 与弹簧的压缩量 $x$ 间的关系如图中实线所示。在另一星球 $N$ 上用完全相同的弹簧，改用物体 $Q$ 完成同样的过程，其 $a$--$x$ 关系如图中虚线所示。假设两星球均为质量均匀分布的球体，已知星球 $M$ 的半径是星球 $N$ 的 3 倍，则\gkblank

\gkfig{figures/q21-ax}

\begin{gkopts}(1)
  \task A．$M$ 与 $N$ 的密度相等
  \task B．$Q$ 的质量是 $P$ 的 3 倍
  \task C．$Q$ 下落过程中的最大动能是 $P$ 的 4 倍
  \task D．$Q$ 下落过程中弹簧的最大压缩量是 $P$ 的 4 倍
\end{gkopts}
\end{question}

\gkpart{三、非选择题}

\textbf{（共 62 分。第 22--25 题为必考题，每个试题考生都必须作答。第 33--34 题为选考题，考生根据要求作答。）}

\gksubpart{（一）必考题（共 47 分）}

\begin{question}
\gkti{22}（5 分）某小组利用打点计时器对物块沿倾斜的长木板加速下滑时的运动进行研究。物块拖动纸带下滑，打出的纸带一部分如图所示。已知打点计时器所用交流电的频率为 \SI{50}{Hz}，纸带上标出的每两个相邻点之间还有 4 个打出的点未画出。在 $A$、$B$、$C$、$D$、$E$ 五个点中，打点计时器最先打出的是\gkblank 点；在打出 $C$ 点时物块的速度大小为\gkblank \si{\metre\per\second}（保留 3 位有效数字）；物块下滑的加速度大小为\gkblank \si{\metre\per\second\squared}（保留 2 位有效数字）。

\gkfig{figures/q22-tape}
\end{question}

\begin{question}
\gkti{23}（10 分）某同学要将一量程为 \SI{250}{\micro\ampere} 的微安表改装为量程为 \SI{20}{\milli\ampere} 的电流表。该同学测得微安表内阻为 \SI{1200}{\ohm}，经计算后将一阻值为 $R$ 的电阻与该微安表连接，进行改装。然后利用一标准毫安表，根据图（a）所示电路对改装后的电表进行检测（虚线框内是改装后的电表）。

\gkfig{figures/q23-circuit}

\begin{enumerate}[leftmargin=2em, label=(\arabic*)]
  \item 根据图（a）和题给条件，将（b）中的实物连接。
  \item 当标准毫安表的示数为 \SI{16.0}{\milli\ampere} 时，微安表的指针位置如图（c）所示，由此可以推测出改装的电表量程不是预期值，而是\gkblank。（填正确答案标号）
        \begin{gkopts}(2)
          \task A．\SI{18}{\milli\ampere}
          \task B．\SI{21}{\milli\ampere}
          \task C．\SI{25}{\milli\ampere}
          \task D．\SI{28}{\milli\ampere}
        \end{gkopts}
  \item 产生上述问题的原因可能是\gkblank。（填正确答案标号）
        \begin{gkopts}(2)
          \task A．微安表内阻测量错误，实际内阻大于 \SI{1200}{\ohm}
          \task B．微安表内阻测量错误，实际内阻小于 \SI{1200}{\ohm}
          \task C．$R$ 值计算错误，接入的电阻偏小
          \task D．$R$ 值计算错误，接入的电阻偏大
        \end{gkopts}
  \item 要达到预期目的，无论测得的内阻值是否正确，都不必重新测量，只需要将阻值为 $R$ 的电阻换为一个阻值为 $kR$ 的电阻即可，其中 $k=$\gkblank。
\end{enumerate}

\gkfig{figures/q23-dial}
\end{question}

\begin{question}
\gkti{24}（12 分）如图，在直角三角形 $OPN$ 区域内存在匀强磁场，磁感应强度大小为 $B$、方向垂直于纸面向外。一带正电的粒子从静止开始经电压 $U$ 加速后，沿平行于 $x$ 轴的方向射入磁场；一段时间后，该粒子在 $OP$ 边上某点以垂直于 $x$ 轴的方向射出。已知 $O$ 点为坐标原点，$N$ 点在 $y$ 轴上，$OP$ 与 $x$ 轴的夹角为 $30^{\circ}$，粒子进入磁场的入射点与离开磁场的出射点之间的距离为 $d$，不计重力。求：
\begin{enumerate}[leftmargin=2em, label=(\arabic*)]
  \item 带电粒子的比荷；
  \item 带电粒子从射入磁场到运动至 $x$ 轴的时间。
\end{enumerate}

\gkfig{figures/q24-opn}
\end{question}

\begin{question}
\gkti{25}（20 分）竖直面内一倾斜轨道与一足够长的水平轨道通过一小段光滑圆弧平滑连接，小物块 $B$ 静止于水平轨道的最左端，如图（a）所示。$t=0$ 时刻，小物块 $A$ 在倾斜轨道上从静止开始下滑，一段时间后与 $B$ 发生弹性碰撞（碰撞时间极短）；当 $A$ 返回到倾斜轨道上的 $P$ 点（图中未标出）时，速度减为 0，此时对其施加一外力，使其在倾斜轨道上保持静止。物块 $A$ 运动的 $v$--$t$ 图像如图（b）所示，图中的 $v_{1}$ 和 $t_{1}$ 均为未知量。已知 $A$ 的质量为 $m$，初始时 $A$ 与 $B$ 的高度差为 $H$，重力加速度大小为 $g$，不计空气阻力。求：
\begin{enumerate}[leftmargin=2em, label=(\arabic*)]
  \item 物块 $B$ 的质量；
  \item 在图（b）所描述的整个运动过程中，物块 $A$ 克服摩擦力所做的功；
  \item 已知两物块与轨道间的动摩擦因数均相等。在物块 $B$ 停止运动后，改变物块与轨道间的动摩擦因数，然后将 $A$ 从 $P$ 点释放，一段时间后 $A$ 刚好能与 $B$ 再次碰上。求改变前后动摩擦因数的比值。
\end{enumerate}

\gkfig{figures/q25-vt}
\end{question}

\gksubpart{（二）选考题（共 15 分。请考生从 2 道物理题中任选一题作答。如果多做，则按所做的第一题计分。）}

\begin{question}
\gkti{33}[物理——选修 3-3]

(1)（5 分）某容器中的空气被光滑活塞封住，容器和活塞绝热性能良好，空气可视为理想气体。初始时容器中空气的温度与外界相同，压强大于外界。现使活塞缓慢移动，直至容器中的空气压强与外界相同。此时，容器中空气的温度\gkblank（填“高于”“低于”或“等于”）外界温度，容器中空气的密度\gkblank（填“大于”“小于”或“等于”）外界空气的密度。

(2)（10 分）热等静压设备广泛用于材料加工中。该设备工作时，先在室温下把惰性气体用压缩机压入到一个预抽真空的炉腔中，然后炉腔升温，利用高温高气压环境对放入炉腔中的材料加工处理，改善其性能。一台热等静压设备的炉腔中某次放入固体材料后剩余的容积为 \SI{0.13}{\metre\cubed}，炉腔抽真空后，在室温下用压缩机将 10 瓶氩气压入到炉腔中。已知每瓶氩气的容积为 \SI{3.2e-2}{\metre\cubed}，使用前瓶中气体压强为 \SI{1.5e7}{\pascal}，使用后瓶中剩余气体压强为 \SI{2.0e6}{\pascal}；室温温度为 $27\,^{\circ}\mathrm{C}$。氩气可视为理想气体。

\begin{enumerate}[leftmargin=2em, label=(\roman*)]
  \item 求压入氩气后炉腔中气体在室温下的压强；
  \item 将压入氩气后的炉腔加热到 $1227\,^{\circ}\mathrm{C}$，求此时炉腔中气体的压强。
\end{enumerate}
\end{question}

\begin{question}
\gkti{34}[物理——选修 3-4]

(1)（5 分）一简谐横波沿 $x$ 轴正方向传播，在 $t=\dfrac{1}{2}T$ 时刻，该波的波形图如图（a）所示，$P$、$Q$ 是介质中的两个质点。图（b）表示介质中某质点的振动图像。下列说法正确的是\gkblank（填正确答案标号。选对 1 个得 2 分，选对 2 个得 4 分，选对 3 个得 5 分；每选错 1 个扣 3 分，最低得分为 0 分）

\gkfig{figures/q34-wave}

\begin{gkopts}(1)
  \task A．质点 $Q$ 的振动图像与图（b）相同
  \task B．在 $t=0$ 时刻，质点 $P$ 的速率比质点 $Q$ 的大
  \task C．在 $t=0$ 时刻，质点 $P$ 的加速度的大小比质点 $Q$ 的大
  \task D．平衡位置在坐标原点的质点的振动图像如图（b）所示
  \task E．在 $t=0$ 时刻，质点 $P$ 与其平衡位置的距离比质点 $Q$ 的大
\end{gkopts}

(2)（10 分）如图，一艘帆船静止在湖面上，帆船的竖直桅杆顶端高出水面 \SI{3}{\metre}。距水面 \SI{4}{\metre} 的湖底 $P$ 点发出的激光束，从水面出射后恰好照射到桅杆顶端，该出射光束与竖直方向的夹角为 $53^{\circ}$（取 $\sin53^{\circ}=0.8$）。已知水的折射率为 $\dfrac{4}{3}$。

\begin{enumerate}[leftmargin=2em, label=(\roman*)]
  \item 求桅杆到 $P$ 点的水平距离；
  \item 船向左行驶一段距离后停止，调整由 $P$ 点发出的激光束方向，当其与竖直方向夹角为 $45^{\circ}$ 时，从水面射出后仍然照射在桅杆顶端，求船行驶的距离。
\end{enumerate}

\gkfig{figures/q34-boat}
\end{question}
